\subsection{معرفی و انگیزه}
این مکانیزم بر مبنای مبانی نظری پیشرفته در ادبیات توجه و فشرده‌سازی \cite{vaswani2017attention} توسعه یافته است.

مکانیزم توجه حالت ماتریسی فاکتورگیری‌شده (\lr{Factorised Matrix State Attention}) با انگیزه جایگزینی ساختار توجه مربعی سنتی با یک مدل حافظه ترتیبی بازگشتی و خطی‌سازی‌شده طراحی شده است. در ترنسفورمرهای استاندارد، هزینه ذخیره‌سازی بافر کلید و مقدار (\lr{KV-Cache}) به صورت خطی با افزایش طول دنباله رشد کرده و نیاز به حفظ کامل تاریخچه توکن‌ها دارد که به بن‌بست پهنای باند حافظه منجر می‌گردد.

ایده محوری ماژول \lr{factorised\_matrix\_state}، به‌کارگیری حافظه مشارکتی ماتریسی (\lr{Associative Matrix Memory}) با پیچیدگی $\mathcal{O}(N)$ است که در آن حالت حافظه در هر گام زمانی به صورت بازگشتی به‌روزرسانی می‌شود. به جای ذخیره و نگهداری یک حالت ماتریسی حجیم و تمام‌رتبه، این ماژول حالت پنهان را به حاصل‌ضرب دو بردار با رتبه تقلیل‌یافته $U V^\top$ تجزیه می‌کند. این رویکرد همراه با گیتینگ زوال نمایی یادگیرنده (\lr{Learnable Decay Gating})، نرمال‌سازی مخرج و آداپتورهای کم‌مرتبه لورا، امکان استنتاج بی‌نهایت کارآمد و انتقال وزن مطمئن را بر روی مدل‌های انکودر و دیکودر فراهم می‌آورد.

\subsection{نتایج تجربی}
ارزیابی ماژول \lr{factorised\_matrix\_state} بر روی بنچمارک GLUE SST-2 با شرایط سخت‌افزاری پایدار بر روی پردازنده گرافیکی NVIDIA RTX 4070 نتایج مندرج در جدول \ref{tab:benchmark_factorised_matrix_state} را به همراه داشته است:

\begin{table}[H]
\centering
\small
\caption{نتایج بنچ‌مارک تجربی ماژول \lr{factorised\_matrix\_state} روی ۴ مدل پایه}
\label{tab:benchmark_factorised_matrix_state}
\begin{tabular}{lccccc}
\toprule
\textbf{مدل پایه \lr{(Base Model)}} & \textbf{دقت \lr{(Accuracy)}} & \textbf{شاخص \lr{F1}} & \textbf{حافظه \lr{(Peak VRAM)}} & \textbf{توان \lr{(Throughput)}} & \textbf{وضعیت} \\
\midrule
\lr{DistilBERT} & \lr{77.06\%} & \lr{0.7608} & \lr{342.49 MB} & \lr{53,181.74 tok/s} & \lr{PASS} \\
\lr{RoBERTa} & \lr{74.77\%} & \lr{0.7727} & \lr{410.26 MB} & \lr{6,820.90 tok/s} & \lr{PASS} \\
\lr{GPT-2} & \lr{70.76\%} & \lr{0.7615} & \lr{349.49 MB} & \lr{8,689.53 tok/s} & \lr{PASS} \\
\lr{Qwen} & \lr{49.43\%} & \lr{--} & \lr{--} & \lr{--} & \lr{PASS} \\
\bottomrule
\end{tabular}
\end{table}

این نتایج مبین پایداری و تعمیم‌پذیری عالی معماری حالت ماتریسی فاکتورگیری‌شده در تمامی مدل‌ها است. دستیابی به دقت ۷۷.۰۶٪ در DistilBERT با توان ۵۳,۱۸۱ توکن بر ثانیه، و دقت‌های بالای ۷۴.۷۷٪ و ۷۰.۷۶٪ به ترتیب در RoBERTa و GPT-2، این مکانیزم را به عنوان یک جایگزین خطی بسیار کارآمد برای لایه‌های توجه مربعی معرفی می‌کند؛ در حالی که مصرف حافظه در تمام مدل‌ها در بازه کم‌نظیر ۳۴۰ تا ۴۱۰ مگابایت محدود مانده است.
